\smallskip

\textbf{Partie A}

\medskip

L'aire de la ZONE 1 est l'aire du triangle OKJ. C'est un triangle rectangle isocèle en J on a donc :

\[\text{aire\,(ZONE 1)} = \text{aire\,(OKJ)} = \dfrac{\mathrm{OJ} \times \mathrm{JK}}{2} = \dfrac{1\times 1}{2} = \boxed{\dfrac{1}{2}}.\]

\smallskip

L'aire de ZONE 2 est l'aire sous la courbe de la fonction carré entre 0 et 1 :

\[\text{aire\,(ZONE 2)} = \int_{0}^{1}x^{2}\d x = \left[\dfrac{x^{3}}{3}\right]_{0}^{1} = \boxed{\dfrac{1}{3}}.\]
\smallskip

$\bullet~~$L'aire de la ZONE 3 est la différence entre l'aire du carré OIJK et la somme des aires de la ZONE 1 et de la ZONE 2 :
\[\text{aire\,(ZONE 3)} = 1  - \left(\dfrac{1}{2} + \dfrac{1}{3}\right) = 1 - \left(\dfrac{3}{6} + \dfrac{2}{6}\right) = \dfrac{6}{6} - \dfrac{5}{6} = \boxed{\dfrac{1}{6}}.\]





\medskip

\textbf{Partie B}

\medskip

\begin{enumerate}
	\item On représente la situation par un arbre pondéré.
	
	\begin{center}
	\psset{treemode=R,nodesep=1.5pt,treesep=1cm,levelsep=2.5cm}
	\pstree{\TR{}}{
	\pstree{\TR{$T$}\naput{$\frac{1}{6}$}}{
		\TR{$G$}\naput{$\frac{1}{2}$}
		\TR{$\overline{G}$}\nbput{$\frac{1}{2}$}}
	\pstree{\TR{$\overline{T}$}\nbput{$\frac{5}{6}$}}{
		\TR{$G$}\naput{$\frac{1}{6}$}
		\TR{$\overline{G}$}\nbput{$\frac{5}{6}$}}}
	\end{center}
	\item Les évènements $T$ et $\overline{T}$ forment une partition de l'univers, donc, d'après la formule des probabilités totales, on a :
	\begin{eqnarray*}
		P(G) &=& P\left(T \cap G\right) + P\left(\overline{T} \cap G\right) = P(T) \times P_{T}(G) + P\left(\overline{T}\right) \times P_{\overline{T}}(G)\\
		&=& \dfrac{1}{6} \times \dfrac{1}{2} + \dfrac{5}{6} \times \dfrac{1}{6} = \dfrac{1}{12} + \dfrac{5}{36} = \dfrac{3}{36} + \dfrac{5}{36}\\
		&=& \dfrac{8}{36} = \dfrac{4 \times 2}{4 \times 9} =  \boxed{\dfrac{2}{9}}.
	\end{eqnarray*}
	\item On veut déterminer la probabilité $P_{G}(T)$ :
	\[P_{G}(T) = \dfrac{P\left(T \cap G\right)}{P(G)} = \dfrac{\dfrac{1}{12}}{\dfrac{2}{9}} = \dfrac{1}{\underset{4}{\cancel{12}}} \times \dfrac{\overset{3}{\cancel{9}}}{2} = \boxed{\dfrac{3}{8}}.\]
\end{enumerate}

\medskip

\textbf{Partie C}

\medskip

\begin{enumerate}
	\item
	\begin{enumerate}
		\item La loi de probabilité de $X_{1}$ est :
		\begin{center}
		\begin{tabularx}{.5\textwidth}{|>{$}c<{$}|*{3}{>{\centering\arraybackslash$}X<{$}|}}
			\hline
			x_{i} & 1 & 2 & 3 \\ \hline
			P\left(X_{1} = x_{i}\right) & \dfrac{1}{2} & \dfrac{1}{3} & \dfrac{1}{6} \rule[-12pt]{0pt}{30pt} \\\hline
		\end{tabularx}
		\end{center}
		On a donc :
		\[E\left(X_{1}\right) = 1 \times \dfrac{1}{2} + 2 \times \dfrac{1}{3} + \cancel{3}\times \dfrac{1}{\underset{2}{\cancel{6}}} = \dfrac{1}{2} + \dfrac{2}{3} + \dfrac{1}{2} = 1 + \dfrac{2}{3} = \dfrac{3}{3} + \dfrac{2}{3} = \boxed{\dfrac{5}{3}}.\]

		\item On calcule la variance en utilisant la formule de König:
		
		\begin{eqnarray*}
			V\left(X_{1}\right) &=& E\left(X_{1}^{2}\right) - \left[E\left(X_{1}\right)\right]^{2}\\
			&=& 1^{2} \times \dfrac{1}{2} + 2^{2} \times \dfrac{1}{3} + 3^{2} \times \dfrac{1}{6} - \left(\dfrac{5}{3}\right)^{2} = \dfrac{1}{2} + \dfrac{4}{3} + \dfrac{9}{6} - \dfrac{25}{9} = \dfrac{9}{18} + \dfrac{24}{18} + \dfrac{27}{18} - \dfrac{50}{18} = \dfrac{10}{18}\\
			&=& \boxed{\dfrac{5}{9}}
		\end{eqnarray*}
	\end{enumerate}
	
	\item
	\begin{enumerate}
		\item L'évènement $\left\{Y = 9\right\}$ correspond à $\left\{X_{1} = 3\right\} \cap \left\{X_{2} = 3\right\} \cap \left\{X_{3} = 3\right\}$. 
		
		Donc, comme les variables $X_1$, $X_2$ et $X_3$ sont indépendantes:
		
		\[P\left(Y = 9\right) = \dfrac{1}{6} \times \dfrac{1}{6} \times \dfrac{1}{6} = \boxed{\dfrac{1}{216}}.\]

		\item D'après la linéarité de l'espérance:
		
		\[E(Y) = E\left(X_{1}\right) + E\left(X_{2}\right) + E\left(X_{3}\right) = \dfrac{5}{3} + \dfrac{5}{3} + \dfrac{5}{3} = \boxed{5}\]
		
		\item Puisque les variables aléatoires $X_{1}$, $X_{2}$ et $X_{3}$ sont indépendantes, on peut utiliser l'additivité de la variance:
		
		\[V(Y) = V\left(X_{1}\right) + V\left(X_{2}\right) + V\left(X_{3}\right) = \dfrac{5}{9} + \dfrac{5}{9} + \dfrac{5}{9} = \dfrac{15}{9} = \dfrac{3 \times 5}{3 \times 3} = \boxed{\dfrac{5}{3}}.\]
	\end{enumerate}
\end{enumerate}


\bigskip
